\section{Definition of STOP}

Within the formal setting established in the preceding section, the outcome denoted as STOP is defined as a diagnostic result concerning the admissibility of continuation under a given enterprise model. The definition introduced here is purely structural and makes no reference to desirability, preference, or decision authority.

An enterprise analysis yields a STOP outcome if, relative to the adopted enterprise model, no continuation state exists that satisfies the model’s internal constraints. In this sense, STOP does not describe an action taken by an organization, but a property of the modeled system: the continuation space defined by the model is empty under the imposed constraints \cite{ashby1956}.

STOP is defined strictly at the level of the model, not at the level of execution or implementation. It is therefore independent of whether actions could be taken in reality outside the model. A STOP outcome asserts only that, within the boundaries of the specified model, no internally coherent evolution can be represented. This distinction between representability within a model and realizability in reality is consistent with established accounts of modeling and anticipatory systems \cite{rosen1985}.

This definition separates STOP from notions of analytical failure. STOP does not indicate missing data, insufficient effort, or incorrect reasoning. It indicates that, given the model and its constraints, continuation is not an admissible analytical outcome. Any attempt to force continuation within the same model would therefore violate internal coherence rather than extend analytical validity.

The definition of STOP is intentionally minimal. It does not specify the nature of the constraints involved, the structure of the state space, or the mechanisms by which admissibility is evaluated. These elements remain abstract at this stage and are addressed only insofar as they are required to establish general properties of STOP as a diagnostic outcome.

Under this definition, STOP is symmetric to continuation: both are possible analytical results, distinguished solely by the structure of the admissible continuation space. STOP arises precisely when this space is exhausted. As such, STOP is neither exceptional nor residual, but a formally well-defined outcome of enterprise analysis within a continuum-based model.
